% GENERATED FILE. Edit canonical lessons, then run npm run generate:books.
% canonical-source-hash: fnv1a64:8d28c8853b33605b
% canonical-lessons: LA-C428-profectio, LA-C428-reditus, LA-C428-peregrinatio, LA-C428-peregrinor, LA-C428-progredior

\chapter{Departure, Return, Travel abroad, To travel abroad, To go forward}
\label{ch:la-departure-return-travel-abroad-to-travel-abroad-to-go-forward}
\hlchaptermodality{428}

\begin{chapteropening}
\textbf{By the end of this chapter:} \emph{I can talk about leaving, coming back and travelling far: a departure, a return, travel abroad, travelling abroad, and going forward.}

Complete the last lesson of chapter 428: I can talk about leaving, coming back and travelling far: a departure, a return, travel abroad, travelling abroad, and going forward.
\end{chapteropening}

\section[profectiō, profectiōnis]{profectiō, profectiōnis — a departure, a setting out}

\label{lesson:LA-C428-profectio}



\begin{quote}
Before the new one: say the Latin for travel money, then the Latin for to set out on.
\end{quote}

\subsection*{You'll want to know: profectiō, profectiōnis}
\textbf{profectiō, profectiōnis} — \textquotedblleft{}a departure, a setting out\textquotedblright{}.

\begin{quote}
Ante profectiōnem mercātor sarcinās parat.
\end{quote}

\begin{quote}
\emph{Before the departure the merchant gets his baggage ready.}
\end{quote}

Third declension, feminine: \emph{profectiō, profectiōnis f.}, the noun of \emph{proficīscor} 'to set out'.

\begin{cousinweb}
\emph{profectiō} is built on \emph{proficīscor} 'to set out' (Lewis \& Short), through its participle \emph{profectus}. EtymDB does not list it.
\end{cousinweb}

\begin{grammarlens}[title={What you've built}]
One more word for telling, explaining and describing at B1.
\end{grammarlens}

\subsection*{Your turn}
\begin{itemize}
\raggedright
  \item \emph{Say it:} \emph{profectiō, profectiōnis}
  \item \emph{Say it:} \emph{profectiō, profectiōnis}, once more
  \item \emph{Say it:} a sentence of your own with \emph{profectiō, profectiōnis}
  \item \emph{Recall:} say the Latin for travel money, then the Latin for to set out on, then say \emph{profectiō, profectiōnis} again
\end{itemize}

\begin{tcolorbox}[breakable,colback=teal!4,colframe=teal!35!black,title={Before you move on}]
What does profectiō, profectiōnis mean? (\textquotedblleft{}A departure, a setting out\textquotedblright{}.) Say it once more.
\end{tcolorbox}
\section[reditus, reditūs]{reditus, reditūs — a return, a coming back}

\label{lesson:LA-C428-reditus}



\begin{quote}
Before the new one: say the Latin for to set out on, then the Latin for a departure.
\end{quote}

\subsection*{You'll want to know: reditus, reditūs}
\textbf{reditus, reditūs} — \textquotedblleft{}a return, a coming back\textquotedblright{}.

\begin{quote}
Uxor nautae reditum marītī exspectat.
\end{quote}

\begin{quote}
\emph{The sailor's wife waits for her husband's return.}
\end{quote}

Fourth declension, masculine: \emph{reditus, reditūs m.} Also income, what comes back from land.

\begin{cousinweb}
\emph{reditus} is built on \emph{redeō} 'to go back' (EtymDB; Lewis \& Short), itself from \emph{red-} 'back' and \emph{eō} 'to go'.
\end{cousinweb}

\begin{grammarlens}[title={What you've built}]
One more word for telling, explaining and describing at B1.
\end{grammarlens}

\subsection*{Your turn}
\begin{itemize}
\raggedright
  \item \emph{Say it:} \emph{reditus, reditūs}
  \item \emph{Say it:} \emph{reditus, reditūs}, once more
  \item \emph{Say it:} a sentence of your own with \emph{reditus, reditūs}
  \item \emph{Recall:} say the Latin for to set out on, then the Latin for a departure, then say \emph{reditus, reditūs} again
\end{itemize}

\begin{tcolorbox}[breakable,colback=teal!4,colframe=teal!35!black,title={Before you move on}]
What does reditus, reditūs mean? (\textquotedblleft{}A return, a coming back\textquotedblright{}.) Say it once more.
\end{tcolorbox}
\section[peregrīnātiō, peregrīnātiōnis]{peregrīnātiō, peregrīnātiōnis — travel abroad, a stay in foreign parts}

\label{lesson:LA-C428-peregrinatio}



\begin{quote}
Before the new one: say the Latin for a departure, then the Latin for a return.
\end{quote}

\subsection*{You'll want to know: peregrīnātiō, peregrīnātiōnis}
\textbf{peregrīnātiō, peregrīnātiōnis} — \textquotedblleft{}travel abroad, a stay in foreign parts\textquotedblright{}.

\begin{quote}
Post longam peregrīnātiōnem senex domum redit.
\end{quote}

\begin{quote}
\emph{After a long time abroad the old man comes home.}
\end{quote}

Third declension, feminine: \emph{peregrīnātiō, peregrīnātiōnis f.} English \emph{peregrination} comes from it.

\begin{cousinweb}
\emph{peregrīnātiō} is built on \emph{peregrīnor} 'to travel abroad' (EtymDB; Lewis \& Short), which comes from \emph{peregrīnus} 'foreign'. English \emph{peregrination} comes from it.
\end{cousinweb}

\begin{grammarlens}[title={What you've built}]
One more word for telling, explaining and describing at B1.
\end{grammarlens}

\subsection*{Your turn}
\begin{itemize}
\raggedright
  \item \emph{Say it:} \emph{peregrīnātiō, peregrīnātiōnis}
  \item \emph{Say it:} \emph{peregrīnātiō, peregrīnātiōnis}, once more
  \item \emph{Say it:} a sentence of your own with \emph{peregrīnātiō, peregrīnātiōnis}
  \item \emph{Recall:} say the Latin for a departure, then the Latin for a return, then say \emph{peregrīnātiō, peregrīnātiōnis} again
\end{itemize}

\begin{tcolorbox}[breakable,colback=teal!4,colframe=teal!35!black,title={Before you move on}]
What does peregrīnātiō, peregrīnātiōnis mean? (\textquotedblleft{}Travel abroad, a stay in foreign parts\textquotedblright{}.) Say it once more.
\end{tcolorbox}
\section[peregrīnor, peregrīnārī]{peregrīnor, peregrīnārī — to travel abroad, to live in foreign parts}

\label{lesson:LA-C428-peregrinor}



\begin{quote}
Before the new one: say the Latin for a return, then the Latin for travel abroad.
\end{quote}

\subsection*{You'll want to know: peregrīnor, peregrīnārī}
\textbf{peregrīnor, peregrīnārī} — \textquotedblleft{}to travel abroad, to live in foreign parts\textquotedblright{}.

\begin{quote}
Mercātor multōs annōs in terrīs longinquīs peregrīnātur.
\end{quote}

\begin{quote}
\emph{The merchant travels abroad in distant lands for many years.}
\end{quote}

Deponent: \emph{peregrīnor, peregrīnārī, peregrīnātus sum}. English \emph{pilgrim} comes from the related adjective \emph{peregrīnus}.

\begin{cousinweb}
\emph{peregrīnor} is built on \emph{peregrīnus} 'foreign' (EtymDB), which goes back to \emph{peregrē} 'abroad', from \emph{per-} 'through, beyond' and \emph{ager} 'land' (Lewis \& Short). English \emph{pilgrim} comes from \emph{peregrīnus}.
\end{cousinweb}

\begin{grammarlens}[title={What you've built}]
One more word for telling, explaining and describing at B1.
\end{grammarlens}

\subsection*{Your turn}
\begin{itemize}
\raggedright
  \item \emph{Say it:} \emph{peregrīnor, peregrīnārī}
  \item \emph{Say it:} \emph{peregrīnor, peregrīnārī}, once more
  \item \emph{Say it:} a sentence of your own with \emph{peregrīnor, peregrīnārī}
  \item \emph{Recall:} say the Latin for a return, then the Latin for travel abroad, then say \emph{peregrīnor, peregrīnārī} again
\end{itemize}

\begin{tcolorbox}[breakable,colback=teal!4,colframe=teal!35!black,title={Before you move on}]
What does peregrīnor, peregrīnārī mean? (\textquotedblleft{}To travel abroad, to live in foreign parts\textquotedblright{}.) Say it once more.
\end{tcolorbox}
\section[prōgredior, prōgredī]{prōgredior, prōgredī — to go forward, to advance}

\label{lesson:LA-C428-progredior}



\begin{quote}
Before the new one: say the Latin for travel abroad, then the Latin for to travel abroad.
\end{quote}

\subsection*{You'll want to know: prōgredior, prōgredī}
\textbf{prōgredior, prōgredī} — \textquotedblleft{}to go forward, to advance\textquotedblright{}.

\begin{quote}
Nox venit, sed viātōrēs ad oppidum prōgrediuntur.
\end{quote}

\begin{quote}
\emph{Night comes, but the travellers go on towards the town.}
\end{quote}

Deponent: \emph{prōgredior, prōgredī, prōgressus sum}. English \emph{progress} comes from its participle.

\begin{cousinweb}
\emph{prōgredior} is built from \emph{prō-} 'forward' and \emph{gradior} 'to step' (EtymDB; Lewis \& Short). English \emph{progress} comes from its participle \emph{prōgressus}.
\end{cousinweb}

\begin{grammarlens}[title={What you've built}]
One more word for telling, explaining and describing at B1.
\end{grammarlens}

\subsection*{Your turn}
\begin{itemize}
\raggedright
  \item \emph{Say it:} \emph{prōgredior, prōgredī}
  \item \emph{Say it:} \emph{prōgredior, prōgredī}, once more
  \item \emph{Say it:} a sentence of your own with \emph{prōgredior, prōgredī}
  \item \emph{Recall:} say the Latin for travel abroad, then the Latin for to travel abroad, then say \emph{prōgredior, prōgredī} again
\end{itemize}

\begin{tcolorbox}[breakable,colback=teal!4,colframe=teal!35!black,title={Before you move on}]
What does prōgredior, prōgredī mean? (\textquotedblleft{}To go forward, to advance\textquotedblright{}.) Say it once more.
\end{tcolorbox}
